\documentclass[11pt,a4paper]{article}
\usepackage[T1]{fontenc}
\usepackage{lmodern,amsmath,amssymb,amsthm,mathtools,microtype}
\usepackage[margin=27mm]{geometry}
\usepackage[colorlinks=true,linkcolor=blue,citecolor=blue,urlcolor=blue]{hyperref}
\newtheorem{theorem}{Theorem}[section]
\newtheorem{proposition}[theorem]{Proposition}
\newtheorem{lemma}[theorem]{Lemma}
\newtheorem{corollary}[theorem]{Corollary}
\theoremstyle{remark}\newtheorem{remark}[theorem]{Remark}
\newcommand{\Z}{\mathbb Z}\newcommand{\Q}{\mathbb Q}
\DeclareMathOperator{\rk}{rank}
\setlength{\emergencystretch}{2em}
\title{Products of Elliptic Indices and Finite Algebraic Correspondences}
\author{}\date{10 November 2026}
\begin{document}\maketitle
\begin{abstract}
Let \(P\) generate a Zariski-dense cyclic subgroup of an abelian
variety \(A\). The triples
\(( [n]P,[h]P,[nh]P )\) are Zariski dense in \(A^3\), even if
\(n\) ranges through an arbitrary infinite set and each \(n\)
is assigned an arbitrary infinite set of values of \(h\).
For elliptic curves we give a direct proof by intersecting
divisors with graphs of multiplication. For general abelian
varieties the assertion follows from Mordell--Lang and a
criterion in terms of affine relations among the indices.
Consequently a fixed finite algebraic correspondence cannot
reconstruct the index product from the two input points
uniformly on these sets. This restriction does not apply to
arbitrary Diophantine definitions or to unrestricted additional
arithmetic conditions.
\end{abstract}

\section{The density theorem}

Throughout, \(k\) is a number field, all Zariski closures are
geometric, and \(A/k\) is a positive-dimensional abelian
variety. Let \(P\in A(k)\), and suppose that
\(\Gamma=\Z P\) is Zariski dense in \(A\).
In particular \(P\) is not torsion.

\begin{theorem}\label{thm:product}
Let \(I\subseteq\Z\) be infinite. For each \(n\in I\), let
\(J_n\subseteq\Z\) be infinite. Then
\[
 \left\{\bigl([n]P,[h]P,[nh]P\bigr):
                 n\in I,\ h\in J_n\right\}
\]
is Zariski dense in \(A^3\).
\end{theorem}

The sets \(J_n\) may depend on \(n\). In particular, arbitrary
fixed congruence restrictions on the auxiliary index are allowed,
provided an infinite set of indices remains for each \(n\).
When \(A\) is an elliptic curve, every nontorsion point satisfies
the density hypothesis on \(\Gamma\).

\subsection{A direct proof for an elliptic curve}

\begin{proof}[Proof of Theorem~\ref{thm:product} when \(A=E\)]
Suppose a proper closed subset \(V\subsetneq E^3\) contains all
the displayed triples. Let \(\mathcal L\) be an ample line bundle
of degree \(d>0\) on \(E\), and put
\[
 \mathcal H=\mathcal L\boxtimes\mathcal L\boxtimes\mathcal L.
\]
A projective embedding given by a sufficiently large power of
\(\mathcal H\), followed by a hypersurface containing \(V\)
but not \(E^3\), supplies a nonzero effective divisor
\[
                      D\in|\mathcal H^{\otimes N}|,
                      \qquad V\subseteq D,
\]
for some fixed \(N>0\).

Only finitely many fibers of the first projection
\(E^3\to E\) are contained in \(D\): any such fiber is an
irreducible divisor and hence an irreducible component of \(D\).
For every other \(n\in I\), the restriction \(D_n\) to the
fiber over \([n]P\) is a nonzero effective divisor on \(E^2\),
with fixed numerical class
\[
                   N(\mathcal L\boxtimes\mathcal L).
\]
The points
\(( [h]P,[nh]P )\), for \(h\in J_n\), are dense in the graph
\[
              G_n=\{(Q,[n]Q):Q\in E\},
\]
because \(J_nP\) is an infinite subset of the curve \(E\).
Thus \(G_n\) is an irreducible component of \(D_n\).

Let \(\mathcal H_0=\mathcal L\boxtimes\mathcal L\) on \(E^2\).
Its self-intersection is \(2d^2\), so
\[
                       \deg_{\mathcal H_0}(D_n)=2Nd^2.
\]
On the other hand,
\[
 \deg_{\mathcal H_0}(G_n)
 =\deg\bigl(\mathcal L\otimes[n]^*\mathcal L\bigr)
 =d(1+n^2),
\]
using \(\deg[n]=n^2\); see~\cite[Chapter~I, Theorem~7.2]{Milne}.
Since \(\mathcal H_0\) is ample and \(D_n-G_n\) is effective,
\[
                         d(1+n^2)\le2Nd^2.
\]
This is impossible for the unbounded set of integers \(n\in I\)
outside the finitely many excluded fibers.
\end{proof}

\section{Affine relations among indices}

We use the following form of the Mordell--Lang theorem:
the intersection of a closed subvariety of an abelian variety
with a finitely generated subgroup is a finite union of cosets
of subgroups of that finitely generated group. This is a theorem,
not an additional conjecture; see~\cite{Faltings}, and the precise
formulation in the introduction of~\cite{GHSZ}.

\begin{proposition}\label{prop:criterion}
Let \(r\ge1\) and \(\Sigma\subseteq\Z^r\).
The set
\[
          \{([m_1]P,\ldots,[m_r]P):\mathbf m\in\Sigma\}
\]
is Zariski dense in \(A^r\) if and only if \(\Sigma\) is not
contained in a finite union of proper rational affine hyperplanes
in \(\Q^r\).
\end{proposition}

\begin{proof}
The product \(\Gamma^r\) is Zariski dense in \(A^r\).
Suppose the displayed set is contained in a proper closed
subvariety \(V\). By Mordell--Lang, \(V\cap\Gamma^r\) is a
finite union of cosets of subgroups of \(\Gamma^r\).
Identify \(\Gamma^r\) with \(\Z^r\), using the nontorsion point
\(P\). None of these subgroups has rank \(r\). Indeed a
rank-\(r\) subgroup has finite index and contains \(m\Gamma^r\)
for some nonzero integer \(m\); this is Zariski dense in \(A^r\),
as multiplication by \(m\) is surjective. Its translate cannot
be contained in \(V\).

Every subgroup of \(\Z^r\) of rank at most \(r-1\) is annihilated
by a nonzero integer linear form. Its cosets are therefore
contained in proper rational affine hyperplanes.
These finitely many hyperplanes contain \(\Sigma\).

Conversely, after clearing denominators, an affine hyperplane
has equation
\[
                   a_1m_1+\cdots+a_rm_r=c,
                   \qquad a_i,c\in\Z,
\]
with some \(a_i\ne0\). The corresponding points lie in the
proper closed subset
\[
             [a_1]Q_1+\cdots+[a_r]Q_r=[c]P
\]
of \(A^r\). This subset is proper because the indicated
homomorphism \(A^r\to A\) is surjective.
A finite union of such proper closed subsets is still proper.
\end{proof}

\begin{proof}[Proof of Theorem~\ref{thm:product} for general \(A\)]
Apply Proposition~\ref{prop:criterion} to the index triples
\((n,h,nh)\).
On these triples, an affine hyperplane becomes
\[
                  (b+cn)h+an=d,
                  \qquad (a,b,c)\ne(0,0,0).
\]
For all but at most one value of \(n\), this equation has at
most one solution \(h\), or none.
Indeed, if \(b+cn\ne0\), it determines \(h\).
If \(b+cn=0\) for all \(n\), then \(b=c=0\), and \(a\ne0\)
makes \(an=d\) possible for at most one \(n\).
Thus finitely many such hyperplanes cannot contain infinitely
many values of \(h\) for every \(n\) in the infinite set \(I\).
\end{proof}

\begin{corollary}\label{cor:polynomial}
Let \(f_1,\ldots,f_r\in\Z[T_1,\ldots,T_s]\), and suppose that
\(1,f_1,\ldots,f_r\) are linearly independent over \(\Q\).
Then
\[
 \{([f_1(\mathbf t)]P,\ldots,[f_r(\mathbf t)]P):
                            \mathbf t\in\Z^s\}
\]
is Zariski dense in \(A^r\).
In particular, for every polynomial \(f\in\Z[T]\) of degree
at least two and every infinite set \(I\subseteq\Z\),
\[
                        \{([n]P,[f(n)]P):n\in I\}
\]
is Zariski dense in \(A^2\).
\end{corollary}

\begin{proof}
The preimage of a proper affine hyperplane is the zero set of a
nonzero polynomial. A finite union of such zero sets cannot
contain \(\Z^s\), since their product is a nonzero polynomial.
For the last statement, the nonzero polynomial
\(an+bf(n)-c\) has only finitely many roots.
Now apply Proposition~\ref{prop:criterion}.
\end{proof}

\section{The finite-correspondence consequence}

\begin{corollary}\label{cor:finite}
There is no proper closed subset \(C\subsetneq A^3\) containing
all the triples in Theorem~\ref{thm:product}.
In particular there is no fixed algebraic correspondence that
has finite fibers over a nonempty open subset of \(A^2\) and
recovers the third coordinate \([nh]P\) from
\(([n]P,[h]P)\) for all these pairs in that open subset.
The same assertion holds for a finite union of such
correspondences, and for finite algebraic auxiliary choices
followed by rational formulas for the third point.
\end{corollary}

\begin{proof}
The first assertion is Theorem~\ref{thm:product}.
For the next assertion, restrict to the nonempty open set where
the correspondence has finite fibers. Each component dominating
the base has dimension at most \(2\dim A\); its image in
\(A^3\) has proper closure. Components whose images do not
dominate the base are contained over a proper closed subset
of the base and also have proper image closure.
The triples of Theorem~\ref{thm:product} remain Zariski dense
after restriction to the inverse image of a nonempty open set
of \(A^2\). They therefore cannot all lie in these finitely
many proper closed images.

Finite algebraic auxiliary choices over the base give a
generically finite variety over that base. Applying rational
formulas and restricting to their domains reduces to the
same dimension argument.
\end{proof}

Here finite fibers means zero-dimensional geometric fibers.
Merely having finitely many rational points in a geometrically
positive-dimensional fiber does not meet this hypothesis.

For an elliptic curve, this excludes replacing the point
\(U=[hn]P\) by any fixed finite algebraic choice constructed
from \(H=[h]P\) and \(Q=[n]P\), while retaining infinitely
many auxiliary indices for every member of an infinite set of
labels. Corollary~\ref{cor:polynomial} similarly excludes a
fixed finite algebraic construction of \([n^2]P\) from
\([n]P\) on any infinite set of indices.

For a restriction of scalars, the general theorem applies once
the chosen cyclic subgroup is known to be Zariski dense in that
abelian variety. Its use does not follow from the cyclic
description of the rational points alone: the Zariski-density
hypothesis must be checked in the relevant variety.

These are restrictions on the indicated reconstruction
strategy. They do not exclude a source that retains only a
specially chosen finite set of auxiliary indices for each
input, changes the arithmetic conditions, uses a different
algebraic group, introduces an additional free rational
parameter, or defines the integers with fewer variables.
The rational-root Pell strategy is not ruled out merely by
these statements: its missing requirement is a complete global
construction of the rational root, together with a sound
argument controlling its denominators.

\begin{thebibliography}{9}
\bibitem{Faltings}
G.~Faltings, \emph{The general case of S.~Lang's conjecture},
in \emph{Barsotti Symposium in Algebraic Geometry},
Perspectives in Mathematics~15, Academic Press, 1994,
pp.~175--182,
\url{https://doi.org/10.1016/B978-0-12-197270-7.50012-7}.
\bibitem{GHSZ}
D.~Ghioca, F.~Hu, T.~Scanlon, and U.~Zannier,
\emph{A variant of the Mordell--Lang conjecture}, 2018,
\url{https://arxiv.org/abs/1804.10561};
author manuscript
\url{https://math.berkeley.edu/~scanlon/papers/Ga_abelian.pdf}.
\bibitem{Milne}
J.~S.~Milne, \emph{Abelian Varieties}, version~2.0, 2008,
\url{https://www.jmilne.org/math/CourseNotes/AV.pdf}.
\end{thebibliography}
\end{document}
